\scanpage{3}
\deleted{\uline{\unsure{Déf|Dèf}}}\quad
\deleted{\uline{Prop} $G\times G\xrightarrow{p}G$ ouvert}

\deleted{\uline{Cor.} Si $P$ un \inserted{\unsure{pr.|ps.}} torseur sous $G$ ($G\times P\xrightarrow[\unsure{1|\vert}]{\sim}P\times P$\inserted{$(g,x)\mapsto(gx,x)$}) alors $\pi:G\times P\longrightarrow P$ ouvert}

\deleted{\unsure{Lemme|lemme}}\quad
\deleted{\uline{Lemme} Soit $X$ schéma sur lequel $G$ opère. Conditions \unsure{équivaltes|équivalentes}}

\deleted{a) $\forall\,S$ sur $k$, et section \deleted{\unsure{un|une}} $a$ de $X_S$ sur $S$, le morphisme $G_S\longrightarrow X_S$, $g\mapsto ga$, est ouvert (resp. ouvert surjectif)}

\deleted{b) Le morphisme \inserted{$\unsure{a_x|a.}:(g,x)\mapsto(gx,x)$} $G\times X\longrightarrow X\times X$ est \deleted{\unsure{univ.}} ouvert\inserted{ univ. sur $X$} (resp. univ. ouvert et surjectif)}

\deleted{c) $\exists$ un $S$ sur $k$, $S\ne\varnothing$, et une \deleted{\unsure{immersion}} section $a$ de $X_S$ sur $S$, tels que}

\uline{Prop} $G$ \inserted{\unsure{schém. en|schém. en.}} groupes sur $S$. $G$ séparé $/S$ ssi section unité est \uline{fermée}

\uline{Cor} $G$ groupe sur corps. Alors $G$ séparé

\uline{Prop} $G$ \unsure{gpe|grp} sur un corps. Alors $p:G\times G\longrightarrow G$ \inserted{univ.} ouvert. Plus \unsure{gén|gén.}, \deleted{\unsure{Cor|Cor.}}\inserted{si} $\deleted{P}X$ \inserted{\unsure{ps|pr}} torseur sous $G$. Alors $G\times\deleted{P}X\xrightarrow{\pi}\deleted{P}X$ est \inserted{univ.} ouvert

\deleted{\uline{Lemme} Soit $G$ opérant sur $X$. Soit \deleted{\illegible{1}} s'il $\exists$ un corps $k'/k$ et $a\in X(k')$ tel que $G_{k'}\longrightarrow X_{k'}$, $g\longmapsto ga$ est ouvert \inserted{univ. $/k'$} (resp. \unsure{univ.} surjectif) c'est vrai pour tt autre p$^{\mathrm{t}}$ de $X$ \unsure{à|:} \unsure{valeurs|valeur} dans un $k''/k$, et \unsure{m.|un} à valeur \unsure{dans|sur} un $S$.}

\deleted{On dit que $X$ est un espace homogène.}

\uline{Prop} \deleted{Supposons que $G\times X\xrightarrow{\pi}X$ soit ouvert (p. ex. $\deleted{P}X$ un \unsure{ps.} torseur),} Alors pour tt $U\subset G$, $V\subset \deleted{P}X$, $U$;$V$ ouverts, \unsure{\ldots}

a) $U.V$ ouvert si $\pi$ ouvert (p ex. $\deleted{P}X$ un ps. tors.)

% slanted
\margin{\unsure{comprendre|comprends} l'\uline{\unsure{évid.}} de c.}
b) $U.V$ stable par $G_0$ \deleted{\unsure{générisation|translation}}

c) $UV$ stable par spécialisation \deleted{\unsure{pour}} (\inserted{si} \deleted{$\pi$ est \unsure{générisant})} \inserted{les $\varphi_{\unsure{x|{}}}:G_{k'}\longrightarrow X_{k'}$ générisants} $x\in X(k')$) \deleted{p. ex. ouverts)}

c$'$) Si $U,V$ \unsure{qu. cpcts|qu. cpts} (et $X$ \inserted{qu} séparé), ou $\unsure{X|k}$ loc. \unsure{noeth|nœth}, $U\times V$ est fermé.

\uline{Cor} Supposons \deleted{\unsure{qu|que}} $\pi:G\times X\longrightarrow X\times X$ ouvert\inserted{ et $X$ quasi-sép.} \inserted{(p. ex $X$ un ps. torseur)}. Alors $X$ est réunion de parties \unsure{qu. cpctes|qu. cptes} \deleted{ouvertes}\inserted{$W_i$} qui sont à la fois ouvertes, fermées, et quasi-compactes [donc stables par $G^0$]. En particulier, les comp. connexes de $X$ sont quasi-compactes.

\margin{?}(Cor.\unsure{?}) Les composantes connexes de $X$ sont les comp irréductib\unsure{les}
% archive sheet number: 2

